% !TeX root = ../../main.tex
\subsection{感嘆文の種類}
    感嘆文とは、驚き、喜び、悲しみ、怒りなどの\uwave{感情や感動をストレートに表現する文}のことである。

    感嘆文の大きな特徴の1つとして、文末に\fitblankbf[]{!}がつくことが挙げられる。

    感嘆文は\fitblankbf[]{what} を使った感嘆文と\fitblankbf[]{how} を使った感嘆文の2パターンしかない。

\subsection{whatを使った感嘆文}

\begin{theorembox}{\textbf{whatを使った感嘆文}}
    \fitblankbf[]{What} + \fitblankbf[]{a [ an ]} + \fitblankbf[]{形容詞} + \fitblankbf[]{名詞} + \fitblankbf[]{!}
\vspace{1.5em}

「なんて \fitblankbf[]{形容詞}な \fitblankbf[]{名詞} でしょう。」

\end{theorembox}
\vspace{1em}

whatを使った感嘆文は\fitblankbf[]{5つ}のブロックから構成されていることを意識しよう!

\begin{enumerate}
    \item What a cool math teacher !
    \dialogueblank{なんてカッコいい数学の先生でしょう。}
    \vspace{0.5em}
    \item なんて大きなケーキでしょう。
    \dialogueblank{What a big cake !}
    \vspace{1em}
    \item \fitblankbf[]{なんてカッコいい数学の先生でしょう。}
    \dialogueblank{What a cool math teacher !}
\end{enumerate}

\subsection{howを使った感嘆文}

\begin{theorembox}{\textbf{howを使った感嘆文}}
    \fitblankbf[]{How} + \fitblankbf[]{形容詞 [ 副詞 ]} + \fitblankbf[]{!}
\vspace{1.5em}

    「なんて \fitblankbf[]{形容詞 [ 副詞 ]} でしょう。」

\end{theorembox}
\vspace{1em}

howを使った感嘆文は\fitblankbf[]{3つ}のブロックから構成されていることを意識しよう!

\begin{enumerate}
    \item How cute !
    \dialogueblank{なんてかわいいのでしょう。}
\vspace{0.5em}
    \item なんて古いのでしょう。
    \dialogueblank{How old !}
\vspace{1em}
    \item \fitblankbf[]{なんてかわいいのでしょう。}
    \dialogueblank{How cute !}
\end{enumerate}

\newpage
